\documentclass[11pt,a4paper]{article}
\usepackage[utf8]{inputenc}
\usepackage[T1]{fontenc}
\usepackage[portuguese]{babel}
\usepackage{xcolor}
\usepackage{hyperref}
\usepackage{geometry}
\usepackage{listings}
\usepackage{tcolorbox}
\usepackage{graphicx}

\geometry{margin=2.5cm}

\definecolor{primary}{RGB}{41, 128, 185}
\definecolor{secondary}{RGB}{44, 62, 80}
\definecolor{codebg}{RGB}{240, 240, 240}
\definecolor{textcol}{RGB}{51, 51, 51}

\hypersetup{
    colorlinks=true,
    linkcolor=primary,
    filecolor=magenta,      
    urlcolor=primary,
}

\lstset{
    backgroundcolor=\color{codebg},
    basicstyle=\ttfamily\small\color{secondary},
    breaklines=true,
    frame=single,
    rulecolor=\color{codebg},
    keywordstyle=\color{primary}\bfseries,
    commentstyle=\color{gray}\itshape,
    stringstyle=\color{teal},
    showstringspaces=false
}

\title{\color{secondary}\Huge\textbf{Muito além do Kubernetes (\textit{Beyond Kubernetes})}}
\author{\color{primary}DevOps com Kubernetes -- 2026}
\date{}

\begin{document}
\maketitle

\section*{\color{primary}O que você aprenderá nesta página}
\begin{itemize}
    \item Opções de plataformas baseadas em Kubernetes
    \item Configurar uma plataforma \textit{serverless} (Knative) e fazer o \textit{deploy} de uma carga de trabalho simples
\end{itemize}

Como o Kubernetes é, na verdade, uma plataforma de base, vamos explorar alguns sistemas e blocos de construção famosos que rodam sobre ele.

\begin{quote}
"Kubernetes é uma plataforma para construir plataformas. É um lugar melhor para se começar, mas não é o objetivo final." — Kelsey Hightower (@kelseyhightower), 27 de Novembro de 2017.
\end{quote}

O OpenShift é um Kubernetes com foco corporativo ("\textit{enterprise}"). Afirmar que você "não usa Kubernetes porque usa OpenShift" seria o equivalente a dizer "Eu não tenho um motor, eu tenho um carro!". Outras opções de Kubernetes prontos para uso em produção incluem o Rancher - que você deve ter visto anteriormente graças ao k3d (k3s) - e o Anthos GKE do Google. Essas são todas opções a serem consideradas quando você precisa tomar a decisão crucial entre qual distribuição de Kubernetes adotar e se deseja usar um serviço gerenciado.

\begin{tcolorbox}[colback=blue!5!white,colframe=blue!75!black,title=Exercício 5.5: Comparação de Plataformas]
Escolha um provedor de serviços (como o Rancher, por exemplo) e compare-o com outro (como o OpenShift). Decida de maneira arbitrária qual provedor de serviço é "melhor" e apresente argumentos em defesa de sua escolha contra a do outro provedor. 

\textit{Para a submissão, uma lista em tópicos (bullet points) é o suficiente.}
\end{tcolorbox}

\section*{\color{primary}Serverless}

O uso de arquiteturas \textit{Serverless} ("sem servidor") ganhou muita popularidade. Quer seja o Google Cloud Run, Knative, OpenFaaS, OpenWhisk, Fission ou Kubeless, todos eles rodam em cima do Kubernetes - ou ao menos são capazes disso. Por isso, uma afirmação como "O Kubernetes está competindo com o serverless" não faz muito sentido.

Para isso, escolheremos o \textbf{Knative}, pois ele é a solução na qual o Google Cloud Run é baseado, apresentando-se como uma opção bastante competente em comparação com outras alternativas \textit{open-source} disponíveis. Isso também nos mantém no tema de "plataformas sobre plataformas", já que ele poderia ser usado para criar sua própria plataforma \textit{serverless}.

O Knative possui seu próprio contrato de execução (\textit{runtime contract}) apoiado pela comunidade. Ele descreve quais tipos de recursos uma aplicação \textbf{deve} e \textbf{deveria} ter para rodar de forma adequada na plataforma. Um requisito essencial é que a própria aplicação precisa ser \textit{stateless} (sem manutenção de estado local) e configurável por meio de variáveis de ambiente. Esse tipo de especificação \textit{open-source} ajuda o projeto a ganhar uma adoção mais ampla.

\begin{tcolorbox}[colback=blue!5!white,colframe=blue!75!black,title=Exercício 5.6: Testando Serverless]
Instale o componente \textit{Knative Serving} no seu cluster k3d.
Para que o Knative funcione localmente no k3d, você precisa criar um cluster \textbf{sem} o Traefik:

\begin{lstlisting}[language=bash]
$ k3d cluster create --port 8082:30080@agent:0 -p 8081:80@loadbalancer --agents 2 --k3s-arg "--disable=traefik@server:0" --image rancher/k3s:v1.34.1-k3s1
\end{lstlisting}

Em seguida, siga os guias de instalação via YAML oficiais. Para a camada de rede, você pode escolher o \textit{Kourier}. Na seção "\textit{Configure DNS}", selecione \textit{Magic DNS (slip.io)}.

Você poderá se deparar com alguns pods em CrashLoopBackOff:
\begin{lstlisting}[language=bash]
$ kubectl get pods -n knative-serving
NAME                                      READY   STATUS             RESTARTS      AGE
activator-67855958d-w2ws8                 0/1     Running            0             64s
webhook-5446675b97-2ngh6                  0/1     CrashLoopBackOff   3 (12s ago)   64s
\end{lstlisting}

Verifique os logs dos \textit{pods} que estão falhando para descobrir como consertar o problema. Note que você pode acessar o serviço localmente da sua máquina host usando:

\begin{lstlisting}[language=bash]
curl -H "Host: hello.default.192.168.240.3.sslip.io" http://localhost:8081
\end{lstlisting}

Onde \textit{Host} é a URL que você obteve ao rodar o comando: \texttt{kubectl get ksvc}
\end{tcolorbox}

\begin{tcolorbox}[colback=blue!5!white,colframe=blue!75!black,title=Exercício 5.7: Faça o deploy para o serverless]
Modifique o serviço Ping-pong do aplicativo "Log Output" para utilizar \textit{serverless}.

DICAS:
\begin{itemize}
    \item Utilize o DNS do serviço totalmente qualificado (\textit{fully qualified}) do Kubernetes ao chamar os serviços. Em vez de usar \texttt{http://pingpong}, você deve usar \texttt{http://pingpong.exercises.svc.cluster.local}.
    \item Você pode definir o "hostname" da função \texttt{url rewrite} se quiser que o endpoint \textit{pingpong} fique disponível via navegador. O \textit{HTTPRoute} pode ficar assim:
\end{itemize}

\begin{lstlisting}[language=bash]
  rules:
    - matches:
        - path:
            type: PathPrefix
            value: /pingpong
      filters:
        - type: URLRewrite
          urlRewrite:
            hostname: pingpong.exercises.192.168.144.3.sslip.io
            path:
              type: ReplacePrefixMatch
              replacePrefixMatch: /
\end{lstlisting}
\end{tcolorbox}

\begin{tcolorbox}[colback=blue!5!white,colframe=blue!75!black,title=Exercício 5.8: O Panorama (Landscape)]
Acesse e analise o documento oficial \textbf{CNCF Cloud Native Landscape}. Circule o logotipo de cada produto ou projeto que você já usou. Em seguida, use uma cor diferente para circular todas as dependências tecnológicas usadas por esses projetos principais. 

Faça uma lista detalhando os cenários e contextos de uso em que você precisou deles. Qualquer software rodado fora do contexto deste curso pode ser simplesmente rotulado como "uso fora do curso".
\end{tcolorbox}

\end{document}
